\chapter{Mapa de la documentación}\label{cap:mapa}

Este manual condensa lo necesario para instalar, desplegar y operar Qadra, y
remite al detalle que ya existe en el repositorio. Este apéndice indica dónde
está cada tema. Todas las rutas son relativas a la raíz del repositorio.

\section{Puntos de entrada}\label{sec:mapa-entrada}

Conviene empezar por \archivo{README.md} (presentación, comandos comunes,
catálogo de \texttt{despacho-cli} y arranque local) y por
\archivo{docs/README.md}, índice por tema de \archivo{docs/}. Las normas de
contribución están en \archivo{CONTRIBUTING.md}. Cada parte del repositorio
tiene además su propio archivo: \archivo{web/README.md} (con
\archivo{web/apps/qadra/README.md} y \archivo{web/apps/ca/README.md}),
\archivo{pki/README.md}, \archivo{latex/README.md},
\archivo{manual-usuario/README.md} y \archivo{manual-tecnico/README.md}.

\section{Documentos de \texttt{docs/} por tema}\label{sec:mapa-docs}

La carpeta \archivo{docs/} contiene un centenar de documentos Markdown en
su primer nivel y tres subcarpetas: \archivo{docs/adr/}, \archivo{docs/ca/} y
\archivo{docs/usability/}. La \cref{tab:mapa-docs} los agrupa por familia, con
el patrón de nombre y algunas rutas representativas; el índice por tema está
en \archivo{docs/README.md}.

\nota{Los informes de \archivo{docs/} conservan la fecha y el alcance de sus
comprobaciones. Su presencia no acredita una ejecución nueva ni convierte una
capacidad prevista en implementada.}

\begingroup
\footnotesize
\renewcommand{\arraystretch}{1.1}
\begin{longtable}{@{}L{3.1cm} L{4.4cm} L{7.9cm}@{}}
  \caption{Familias de documentos de \texttt{docs/}.}\label{tab:mapa-docs} \\
  \toprule
  \textbf{Tema} & \textbf{Patrón} & \textbf{Rutas representativas} \\
  \midrule
  \endfirsthead
  \toprule
  \textbf{Tema} & \textbf{Patrón} & \textbf{Rutas representativas} \\
  \midrule
  \endhead
  \bottomrule
  \endlastfoot
  Estado y alcance
    & \archivo{docs/technical-closure.md} y afines
    & \archivo{docs/technical-closure.md},
      \archivo{docs/product-completion.md},
      \archivo{docs/production-acceptance.md},
      \archivo{docs/user-guide.md} \\
  Despliegue (\cref{cap:despliegue})
    & \archivo{docs/deployment*.md}
    & \archivo{docs/deployment.md},
      \archivo{docs/deployment-controller-installation.md},
      \archivo{docs/deployment-controller-launcher.md},
      \archivo{docs/deployment-controller-publication.md},
      \archivo{docs/deployment-readiness-deadline.md} \\
  Base de datos, respaldo y restauración (\cref{cap:operacion})
    & \archivo{docs/database-*.md}, \archivo{docs/deployment-restore-*.md}
    & \archivo{docs/database-operations.md},
      \archivo{docs/deployment-backups.md},
      \archivo{docs/backup-capture-manifest.md},
      \archivo{docs/database-restore-commands.md},
      \archivo{docs/deployment-restore-quiesce.md},
      \archivo{docs/restore-process-observation.md} \\
  Autoridad certificadora (\cref{cap:ca})
    & \archivo{docs/ca/*.md}
    & \archivo{docs/ca/README.md},
      \archivo{docs/ca/operations.md},
      \archivo{docs/ca/api.md},
      \archivo{docs/ca/owner-activation.md},
      \archivo{docs/ca/restore.md},
      \archivo{docs/ca/verification.md},
      \archivo{docs/deployment-crl.md} \\
  Contratos HTTP (\cref{cap:api})
    & \archivo{docs/http-*.md}, \archivo{docs/*-api.md}
    & \archivo{docs/http-api.md},
      \archivo{docs/members-api.md},
      \archivo{docs/case-administration-api.md},
      \archivo{docs/document-upload-admission-api.md},
      \archivo{docs/hearings-api.md},
      \archivo{docs/deadlines-api.md} \\
  Identidad, sesión y acceso
    & \archivo{docs/member-*.md}, \archivo{docs/password-reset-*.md},
      \archivo{docs/owner-certificate-*.md}
    & \archivo{docs/member-enrollment-recovery.md},
      \archivo{docs/password-reset-internal.md},
      \archivo{docs/session-editor-recovery.md},
      \archivo{docs/mfa-access-logs.md},
      \archivo{docs/owner-certificate-login-operations.md} \\
  Firma y criptografía en el navegador
    & \archivo{docs/*-signing*.md}
    & \archivo{docs/individual-signing-guide.md},
      \archivo{docs/http-individual-signing.md},
      \archivo{docs/http-contract-signing.md},
      \archivo{docs/browser-local-cryptography.md} \\
  Documentos y formatos
    & \archivo{docs/document-*.md}
    & \archivo{docs/document-format-operations.md},
      \archivo{docs/media-decoder-setup.md},
      \archivo{docs/document-content-api.md} \\
  Audiencias, medidas y recursos
    & \archivo{docs/hearing*.md}, \archivo{docs/precautionary-*.md},
      \archivo{docs/measure-*.md}, \archivo{docs/resource-*.md}
    & \archivo{docs/hearing-results-api.md},
      \archivo{docs/precautionary-hearings-scope.md},
      \archivo{docs/measure-records-api.md},
      \archivo{docs/procedural-resources-scope.md},
      \archivo{docs/resource-hearings.md},
      \archivo{docs/agenda-api.md}, \archivo{docs/alerts-api.md} \\
  Hechos procesales y calendarios
    & \archivo{docs/procedural-*.md}
    & \archivo{docs/procedural-facts.md},
      \archivo{docs/procedural-facts-canonical.md},
      \archivo{docs/procedural-time.md},
      \archivo{docs/judicial-calendars-api.md} \\
  Plazos
    & \archivo{docs/deadline*.md},
      \archivo{docs/legal-corpus-*.md}
    & \archivo{docs/deadline-lifecycle.md},
      \archivo{docs/deadline-arithmetic.md},
      \archivo{docs/deadline-dispatch.md},
      \archivo{docs/deadline-worker.md},
      \archivo{docs/four-front-closure.md} \\
  Verificación y pruebas (\cref{cap:mantenimiento})
    & \archivo{docs/*verification*.md}, \archivo{docs/access-audit-*.md}
    & \archivo{docs/verification-report.md},
      \archivo{docs/academic-report-verification.md},
      \archivo{docs/access-audit-matrix.md},
      \archivo{docs/browser-preparation-checks.md},
      \archivo{docs/ci-runner-operations.md} \\
  Usabilidad
    & \archivo{docs/usability/}
    & \archivo{docs/usability/README.md},
      \archivo{docs/usability/protocol.md},
      \archivo{docs/usability/tasks.md},
      \archivo{docs/usability/facilitator.md} \\
\end{longtable}
\endgroup

\section{Anexos del reporte técnico}\label{sec:mapa-anexos}

El reporte técnico (\archivo{latex/main.tex}) incluye cinco anexos que
amplían varios capítulos de este manual. Sus fuentes están en
\archivo{latex/chapters/}.

\begin{table}[H]
  \centering
  \caption{Anexos del reporte técnico.}\label{tab:mapa-anexos}
  \begin{tabularx}{\textwidth}{C{1.2cm} Y L{5.4cm} L{2.2cm}}
    \toprule
    \textbf{Anexo} & \textbf{Título} & \textbf{Fuente} & \textbf{Véase} \\
    \midrule
    A & Catálogo de formatos canónicos y migraciones
      & \archivo{anexos.tex} & \cref{cap:arquitectura} \\
    B & Manual de la interfaz de línea de comandos
      & \archivo{anexo-b-cli.tex} & \cref{cap:cli} \\
    C & Guía de verificación independiente con OpenSSL
      & \archivo{anexo-c-verificacion.tex} & \cref{cap:verificacion} \\
    D & Scripts de la autoridad certificadora interna
      & \archivo{anexo-d-pki.tex} & \cref{cap:ca} \\
    E & Matriz de pruebas y evidencia de ejecución
      & \archivo{anexo-e-matriz.tex} & \cref{cap:mantenimiento} \\
    \bottomrule
  \end{tabularx}
\end{table}

Los anexos B, D y E se componen de varios archivos: el archivo indicado en la
tabla incluye los demás de su misma letra (\archivo{anexo-b-*.tex},
\archivo{anexo-d-*.tex} y \archivo{anexo-e-*.tex}).

\section{Manual de usuario}\label{sec:mapa-usuario}

El manual de usuario es un documento independiente en
\archivo{manual-usuario/}, con la misma estructura de archivos que este
manual. Describe el uso de la aplicación, no su operación, y se organiza en
una introducción y tres partes: A, \emph{Uso del sistema en el despacho};
B, \emph{Cliente y verificación de evidencia}, y C, \emph{Referencia rápida}
(mensajes del sistema y solución de problemas desde la interfaz).

Las capturas de pantalla previstas se relacionan en
\archivo{manual-usuario/capturas.md}.

\section{Decisiones de arquitectura}\label{sec:mapa-adr}

La carpeta \archivo{docs/adr/} contiene los registros de decisión de
arquitectura (ADR), con nombre \archivo{NNNN-titulo.md}; dos registros
comparten el número 0048. La
\cref{tab:mapa-adr} recoge los que más inciden en la instalación, el
despliegue y la operación; el título se da traducido y el archivo se
encuentra en \archivo{docs/adr/}.

\begingroup
\footnotesize
\renewcommand{\arraystretch}{1.1}
\begin{longtable}{@{}C{1.1cm} L{6.6cm} L{7.7cm}@{}}
  \caption{ADR más relevantes para operar el sistema.}\label{tab:mapa-adr} \\
  \toprule
  \textbf{Núm.} & \textbf{Título} & \textbf{Archivo en \texttt{docs/adr/}} \\
  \midrule
  \endfirsthead
  \toprule
  \textbf{Núm.} & \textbf{Título} & \textbf{Archivo en \texttt{docs/adr/}} \\
  \midrule
  \endhead
  \bottomrule
  \endlastfoot
  0001 & Versión mínima de Rust y fijación de dependencias
    & \archivo{0001-msrv-and-dependency-pinning.md} \\
  0006 & Formato del paquete de evidencia
    & \archivo{0006-evidence-package-format.md} \\
  0007 & Anclaje de la cadena de auditoría
    & \archivo{0007-audit-chain-anchoring.md} \\
  0009 & Autoridad local de sellado de tiempo
    & \archivo{0009-local-timestamp-authority.md} \\
  0012 & Sesiones revocables y control de acceso por roles
    & \archivo{0012-revocable-sessions-and-rbac.md} \\
  0044 & Costo de referencia del resumen de contraseñas
    & \archivo{0044-reference-password-hashing-cost.md} \\
  0048 & Despliegue privado y versionado sobre SSH
    & \archivo{0048-private-versioned-deployment.md} \\
  0048 & Coordinación de migraciones acotada a un esquema de PostgreSQL
    & \archivo{0048-schema-scoped-migration-lock.md} \\
  0062 & Validación acotada de PostgreSQL al construir el servidor
    & \archivo{0062-scoped-postgres-startup-validation.md} \\
  0063 & Mantenimiento recuperable de la CRL en un despliegue privado
    & \archivo{0063-recoverable-crl-maintenance.md} \\
  0064 & Vida absoluta de la sesión y actividad explícita
    & \archivo{0064-explicit-session-activity.md} \\
  0078 & Operador de CA separado y operaciones de certificado recuperables
    & \archivo{0078-separate-ca-operator-and-recoverable-operations.md} \\
  0079 & Claves locales en el navegador y firmas sobre bytes exactos
    & \archivo{0079-browser-local-cryptography.md} \\
  0080 & Renovación programada de la CRL cerca de su vencimiento
    & \archivo{0080-scheduled-crl-renewal.md} \\
  0085 & Validación de evidencia en el navegador y manifiesto sellado
    & \archivo{0085-in-browser-evidence-validation.md} \\
  0086 & Firma bilateral de contratos
    & \archivo{0086-bilateral-contract-signing.md} \\
\end{longtable}
\endgroup

Las decisiones sobre la integración continua (entre otras, la 0027, de la
0045 a la 0055 salvo las dos 0048, la 0081, la 0082 y la 0088) sólo interesan
a quien mantenga los ejecutores de CI; su operación se describe en
\archivo{docs/ci-runner-operations.md}.
